% =====================================================================
% 5. TARTIŞMA (Discussion) — EN KRİTİK BÖLÜM (%25)
% Evidence base: docs/results_progress.md (key findings), docs/experiment_log.md,
%                docs/report/AGENT_REPORT_RULES.md §4-§7.
% Address the assignment §5.5 questions explicitly. Report neutral/negative
% results honestly. No SOTA, no protocol-mismatched head-to-head claims.
% =====================================================================
\section{Tartışma}
\label{sec:discussion}

% TODO: Hangi omurga ağ daha iyi özellik öğrendi ve neden?
%       Frozen'da EfficientNetB0 en iyi tekil (F1 0.5586); CV'de weighted füzyon kazandı.
\TODO{Hangi CNN daha iyi özellik öğrendi? (frozen E en iyi tekil; gerekçe).}

% TODO: Füzyon performansı artırdı mı? Concat vs weighted CI ayrımı
%       (exp 11 lower 0.5814 > exp 10 upper 0.5799). GMU concat'i geçemedi (Δ≈0.002).
\TODO{Füzyon yardımcı oldu mu? concat vs weighted CI ayrımı; GMU nötr sonucu.}

% TODO: En iyi kombinasyon — üçlü weighted fine-tune (exp 11), CV F1 0.5892,
%       pooled 0.6000. Fold-0 sıralamasının yanıltıcı olduğu (concat>weighted) ve
%       5-fold ile tersine döndüğü açıklanacak (metodolojik ders).
\TODO{En iyi kombinasyon (exp 11) ve fold-0 yanılgısının açıklaması.}

% TODO: Güçlü/zayıf yönler — nadir sınıflar (support ≤ 10) F1=0; hata matrisi
%       rastgele değil, anlamsal/görsel olarak yakın çoğunluk sınıfına soğuruluyor.
\TODO{Güçlü/zayıf yönler; nadir-sınıf çöküşü (support ≤ 10) ve hata yapısı.}

% TODO: Sınıf dengesizliği etkisi — accuracy 0.87 vs macro-F1 0.60 farkı;
%       weighted-F1 0.87 / MCC 0.86. macro-F1'in neden headline olduğu.
\TODO{Dengesizlik etkisi; accuracy$\leftrightarrow$macro-F1 farkı; macro-F1 gerekçesi.}

% TODO: Hangi transfer learning yaklaşımı daha başarılı ve NE KADAR SÜRDÜ?
%       Fine-tune CV ortalamasında frozen'ı geçti. Süre: Tablo 4'e bağlı (HENÜZ YOK).
\TODO{Frozen vs fine-tune sonucu + süre yorumu (Tablo 4 üretilince bağlanacak).}

% TODO: Literatür ile dürüst konumlandırma — bizim ~0.60 macro-F1, Borgli'nin
%       resmi-split ~0.62 bandında; Effimix 0.97 farklı protokol (augment+balance+
%       non-official split) → doğrudan kıyaslanamaz. Saklama, açıkla (§8 dürüstlük).
\TODO{Dürüst literatür konumlandırması (protokol farkı; SOTA iddiası YOK).}

% NOTE: En önemli sonuçların yorumu ve sınırlamalar öğrenci tarafından elle
%       gözden geçirilecek (AGENT_REPORT_RULES.md §9).
